% Comment 5: verified against normalization, split and graph code and frozen artifacts.
\subsection{\rev{repoDB normalization and cross-resource entity alignment}}
\label{sec:repodb_alignment}
\begin{revisionblock}
\textbf{Normalization and target definition.}
The repoDB preparation code cleans HTML and whitespace, converts identifier strings to uppercase, and extracts DrugBank IDs matching \texttt{DB} followed by five digits and UMLS CUIs matching \texttt{C} followed by seven digits. It reads the structured drug and indication fields, falling back to the corresponding raw text fields when necessary. All 10,800 input rows have valid identifiers. Drugs are represented as \texttt{Compound::DBxxxxx} and diseases as \texttt{Disease::UMLS:Cxxxxxxx}; names are retained for display rather than used as cross-resource join keys. Approved records define positive targets. Terminated, withdrawn, suspended, or failed records form a separate diagnostic relation, not positive indication labels or assumed mechanistic contraindications. Deduplication by canonical drug--disease pair gives 6,677 approved pairs and 3,281 failed-like pairs. The 21 pairs appearing in both groups retain their approved label and are removed from the diagnostic group, leaving 3,260 diagnostic pairs.

\textbf{Alignment scope and unmatched entities.}
The join uses exact canonical DrugBank IDs in the prepared DRKG entity map, which contains 13,918 entities; it is not a mapping against the complete raw DRKG. Across approved and diagnostic pairs, 1,130 of 1,571 unique compounds (71.93\%) match this map. Among the 1,519 approved compounds, 1,106 match (72.81\%). For each of the 2,051 distinct disease CUIs, the code probes three literal forms: \texttt{Disease::UMLS:Cxxxxxxx}, \texttt{Disease::Cxxxxxxx}, and the bare CUI. None matches. These checks test identifier equality only; they do not establish the absence of semantic correspondences. No ontology crosswalk, synonym-based disease alignment, or disease-to-gene identifier conversion is performed. Diseases remain local UMLS nodes, and unmatched compounds are retained rather than excluded from the benchmark. Thus local graph-ID coverage and cross-resource alignment coverage are different quantities.

\textbf{Reuse of DRKG evidence.}
The graph starts with training approved-indication edges and conflict-cleaned diagnostic edges. It then scans the prepared DRKG training graph for compound--gene edges incident to repoDB compounds with matching IDs, preserving the original relation, direction, and gene ID. Gene--gene bridges are retained only when both endpoints belong to the gene set touched by those compound--gene edges. No DRKG disease--gene edges are imported through a UMLS disease mapping. Diagnostic edges are capped at 80 per disease and 80 per compound; compound--gene edges at 80 per compound; gene--gene edges at 40 per gene and 150,000 overall. Edges are deduplicated, and the caps are applied in the stored input order. This yields 5,677 training approved edges, 3,222 diagnostic edges, 39,017 compound--gene edges, and 42,888 gene--gene edges: 90,804 triples over 8,818 entities and 63 relations. The drop from 3,260 to 3,222 diagnostic pairs occurs during capped graph construction. All task entities are included in the local ID maps even if they lack imported DRKG evidence.

\textbf{Splitting and leakage checks.}
The 6,677 unique approved triples are split into 5,677/500/500 training/validation/test triples with seed 2028. A triple is eligible for holdout only if its compound and disease each retain another approved training edge. The resulting 1,519 training compounds define the candidate universe. The splits are disjoint, every validation/test compound and disease remains in the training target-relation universe, and no held-out approved triple occurs in the constructed evidence graph. Diagnostic edges for held-out positive pairs are also excluded. These are exact-triple and entity-coverage checks, not claims that all semantic overlap between resources has been removed. Evaluation candidate lists are produced without gold injection; canonical and numeric IDs are retained separately from display names in the prepared graph--LLM inputs.

The alignment therefore supplies drug-side molecular evidence and gene bridges while leaving disease-side mechanistic context limited to the local graph. This defines the scope of the repoDB setting: evaluation under partial cross-resource alignment, rather than a fully aligned UMLS--DRKG benchmark.
\end{revisionblock}
